% sec1_7_part1.tex — Section 1.7 (part 1): Electricity Supply Application (pp. 60--65)

Nerlove's 1963 paper is a classic study of returns to scale in a regulated industry.
It illustrates the techniques of this chapter and presents a few more not yet covered.

\subsection*{The Electricity Supply Industry}

At the time of Nerlove's writing, the U.S.\ electric power supply industry had the
following features:
\begin{enumerate}
\item Privately owned local monopolies supply power on demand.
\item Rates (electricity prices) are set by the utility commission.
\item Factor prices (e.g., the wage rate) are given to the firm, either because of
  perfect competition in the market for factor inputs or through long-term contracts
  with labor unions.
\end{enumerate}
These institutional features will be relevant when we examine whether the OLS is an
appropriate estimation procedure.\footnote{Thanks to the deregulation of the industry
since the time of Nerlove's writing, multiple firms are now allowed to compete in the
same local market, and the strict price control has been lifted in many states.  So the
first two features no longer characterize the industry.}

\subsection*{The Data}

Nerlove assembled a cross-section data set on 145 firms in 44 states in the year 1955
for which data on all the relevant variables were available.  The variables in the data
are total costs, factor prices (the wage rate, the price of fuel, and the rental price
of capital), and output.  Although firms own capital (such as power plants, equipment,
and structures), the standard investment theory of Jorgenson (1963) tells us that the
firm should behave as if it rents capital at a rental price called the ``user cost of
capital,'' which is defined as $(r + \delta) \cdot p_I$, where $r$ here is the real
interest rate, $\delta$ is the depreciation rate, and $p_I$ is the price of capital goods.
For this reason capital input can be treated as if it is a variable factor of production,
just like labor and fuel inputs.

\subsection*{Why Do We Need Econometrics?}

Why do we need a fancy econometric technique like OLS to determine returns to scale?
The reason is that each firm can have a different AC curve.  If firms face different
factor prices, then the average cost is less for firms facing lower factor prices.
The effect of factor prices on the AC curve has to be isolated somehow.  The approach
taken by Nerlove, which became a standard econometric practice, is to estimate a
parameterized cost function.

Another factor that shifts the individual AC curve is the level of production efficiency.
If more efficient firms produce more output, then it is possible that the individual AC
curve is upward sloping but the line connecting the observed combination of the average
cost and output is downward sloping.  To illustrate, consider a competitive industry
described in Figure~1.6, where the AC and MC (marginal cost) curves are drawn for two
firms competing in the same market.

\begin{figure}[htbp]
\centering
\includegraphics[width=0.9\textwidth,trim=50 400 50 400,clip]{figures/fig_1_6.png}
\caption{Output Determination}
\label{fig:1.6}
\end{figure}

\subsection*{The Cobb-Douglas Technology}

To derive a parameterized cost function, we start with the Cobb-Douglas production
function
\begin{equation}
  Q_i = A_i \, x_{i1}^{\alpha_1} \, x_{i2}^{\alpha_2} \, x_{i3}^{\alpha_3},
  \label{eq:1.7.1}
\end{equation}
where $Q_i$ is firm $i$'s output, $x_{i1}$ is labor input for firm $i$, $x_{i2}$ is
capital input, and $x_{i3}$ is fuel.  $A_i$ captures unobservable differences in
production efficiency (this term is often called \textbf{firm heterogeneity}).  The sum
$\alpha_1 + \alpha_2 + \alpha_3 \equiv r$ is the degree of returns to scale.  Thus, it is
assumed \emph{a priori} that the degree of returns to scale is constant (this should not
be confused with constant returns to scale, which is that $r = 1$).

We know from microeconomics that the cost function associated with the Cobb-Douglas
production function is Cobb-Douglas:
\begin{equation}
  TC_i = r \cdot (A_i \, \alpha_1^{\alpha_1} \, \alpha_2^{\alpha_2} \,
    \alpha_3^{\alpha_3})^{-1/r} \, Q_i^{1/r} \, p_{i1}^{\alpha_1/r} \,
    p_{i2}^{\alpha_2/r} \, p_{i3}^{\alpha_3/r},
  \label{eq:1.7.2}
\end{equation}
where $TC_i$ is total costs for firm $i$.  Taking logs, we obtain the following
log-linear relationship:
\begin{equation}
  \log(TC_i) = \mu_i + \frac{1}{r}\log(Q_i) + \frac{\alpha_1}{r}\log(p_{i1})
    + \frac{\alpha_2}{r}\log(p_{i2}) + \frac{\alpha_3}{r}\log(p_{i3}),
  \label{eq:1.7.3}
\end{equation}
where $\mu_i = \log[r \cdot (A_i \, \alpha_1^{\alpha_1} \, \alpha_2^{\alpha_2} \,
\alpha_3^{\alpha_3})^{-1/r}]$.  The equation is said to be \textbf{log-linear} because
both the dependent variable and the regressors are logs.  Coefficients in log-linear
equations are \textbf{elasticities}.  The degree of returns to scale, which
in~\eqref{eq:1.7.3} is the reciprocal of the output elasticity of total costs, is
independent of the level of output.

Now let $\mu \equiv \E(\mu_i)$ and define $\varepsilon_i \equiv \mu_i - \mu$ so that
$\E(\varepsilon_i) = 0$.  This $\varepsilon_i$ represents the inverse of the firm's
production efficiency relative to the industry's average efficiency; firms with positive
$\varepsilon_i$ are high-cost firms.  With this notation, \eqref{eq:1.7.3} becomes
\begin{equation}
  \log(TC_i) = \beta_1 + \beta_2 \log(Q_i) + \beta_3 \log(p_{i1})
    + \beta_4 \log(p_{i2}) + \beta_5 \log(p_{i3}) + \varepsilon_i,
  \label{eq:1.7.4}
\end{equation}
where
\begin{equation}
  \beta_1 = \mu, \quad \beta_2 = \frac{1}{r}, \quad \beta_3 = \frac{\alpha_1}{r},
  \quad \beta_4 = \frac{\alpha_2}{r}, \quad \text{and} \quad
  \beta_5 = \frac{\alpha_3}{r}.
  \label{eq:1.7.5}
\end{equation}
Thus, the cost function has been cast in the regression format of Assumption~1.1
with $K = 5$.

\subsection*{How Do We Know Things Are Cobb-Douglas?}

The Cobb-Douglas functional form is certainly a very convenient parameterization of
technology.  But how do we know that the true production function is Cobb-Douglas?
The Cobb-Douglas form satisfies the properties, such as diminishing marginal
productivities, that we normally require for the production function, but it is
certainly not the only functional form with those desirable properties.  A number of
more general functional forms have been proposed in the literature, but the
Cobb-Douglas form, despite its simplicity, has proved to be a surprisingly good
description of technology.

\subsection*{Are the OLS Assumptions Satisfied?}

To justify the use of least squares, we need to make sure that Assumptions 1.1--1.4
are satisfied for the equation~\eqref{eq:1.7.4}.  Evidently, Assumption~1.1 (linearity)
is satisfied with
\[
  y_i = \log(TC_i), \quad
  \vec{x}_i = (1,\, \log(Q_i),\, \log(p_{i1}),\, \log(p_{i2}),\, \log(p_{i3}))'.
\]
There is no reason to expect that the regressors in~\eqref{eq:1.7.4} are perfectly
multicollinear.  Indeed, in Nerlove's data set, $\rank(\vec{X}) = 5$ and $n = 145$,
so Assumption~1.3 (no multicollinearity) is satisfied as well.

In verifying strict exogeneity (Assumption~1.2), the features of the electricity
industry are relevant.  It is reasonable to assume, as in most cross-section data, that
$\vec{x}_i$ is independent of $\varepsilon_j$ for $i \neq j$.  So the question is
whether $\vec{x}_i$ is independent of $\varepsilon_i$.  According to the third feature
of the industry, factor prices are given to the firm with no regard for the firm's
efficiency, so it is eminently reasonable to assume that factor prices are independent
of $\varepsilon_i$.

What about output?  Since the firm's output is supplied on demand (the first feature of
the industry), output depends on the price of electricity set by the utility commission
(the second feature).  If the regulatory scheme is such that the price is determined
regardless of the firm's efficiency, then $\log(Q_i)$ and $\varepsilon_i$ are
independently distributed.  On the other hand, if the price is set to cover the average
cost, then the firm's efficiency affects output and output in this case is
\textbf{endogenous}, being correlated with the error term.

There is no \emph{a priori} reason to suppose that homoskedasticity is satisfied.
Indeed, the plot of residuals to be shown shortly suggests a failure of this condition.

\subsection*{Restricted Least Squares}

The equation~\eqref{eq:1.7.4} is \textbf{overidentified} in that its five coefficients,
being functions of the four technology parameters (which are $\alpha_1$, $\alpha_2$,
$\alpha_3$, and $\mu$), are not free parameters.  We can easily see that
from~\eqref{eq:1.7.5}: $\beta_3 + \beta_4 + \beta_5 = 1$ (recall:
$r \equiv \alpha_1 + \alpha_2 + \alpha_3$).  This is a reflection of the generic
property of the cost function that it is linearly homogeneous in factor prices.  Indeed,
multiplying total costs $TC_i$ and all factor prices $(p_{i1}, p_{i2}, p_{i3})$ by a
common factor leaves the cost function~\eqref{eq:1.7.4} intact if and only if
$\beta_3 + \beta_4 + \beta_5 = 1$.

Estimating the equation by least squares while imposing \emph{a priori} restrictions on
the coefficient vector is the restricted least squares.  In the present example, to
impose the homogeneity restriction $\beta_3 + \beta_4 + \beta_5 = 1$ on the cost
function, we take any one of the factor prices, say $p_{i3}$, and subtract
$\log(p_{i3})$ from both sides of~\eqref{eq:1.7.4} to obtain
\begin{equation}
  \log\!\left(\frac{TC_i}{p_{i3}}\right)
  = \beta_1 + \beta_2 \log(Q_i) + \beta_3 \log\!\left(\frac{p_{i1}}{p_{i3}}\right)
    + \beta_4 \log\!\left(\frac{p_{i2}}{p_{i3}}\right) + \varepsilon_i.
  \label{eq:1.7.6}
\end{equation}
There are now four coefficients in the regression, from which unique values of the four
technology parameters can be determined.  The restricted least squares estimate of
$(\beta_1, \ldots, \beta_4)$ is simply the OLS estimate of the coefficients
in~\eqref{eq:1.7.6}.  The restricted least squares estimate of $\beta_5$ is the value
implied by the estimate of $(\beta_1, \ldots, \beta_4)$ and the restriction.
